% Source: https://artofproblemsolving.com/community/c3h1937899_maths_olympiad
% Stable identity: aops-topic-1937899-q01
The conditions give $AD=\dfrac{1}{4}AB$, $BE=\dfrac{1}{4}BC$, $CF=\dfrac{1}{4}CA$, and so $BD=\dfrac{3}{4}AB$, $CE=\dfrac{3}{4}CB$, $AF=\dfrac{3}{4}AC$.

\begin{asy}
usepackage("amsmath");
size(6cm);
defaultpen(fontsize(10pt));
pair A=(1.2,4),B=(0,0),C=(6,0);
pair D=A+(B-A)/4,Ee=B+(C-B)/4,F=C+(A-C)/4;
draw(A--B--C--cycle,linewidth(1));
filldraw(D--Ee--F--cycle,palegray,linewidth(0.9));
dot(D);dot(Ee);dot(F);
label("$A$",A,N);label("$B$",B,SW);label("$C$",C,SE);
label("$D$",D,W);label("$E$",Ee,S);label("$F$",F,NE);
label("\mbox{}",(min(currentpicture,user=true)+max(currentpicture,user=true))/2,black); // keeps the text after the figure black when LaTeX typesets the labels inline
\end{asy}

Two triangles sharing an angle have areas in the ratio of the products of the sides enclosing it. Hence
\[
\dfrac{[ADF]}{[ABC]}=\dfrac{AD}{AB}\cdot\dfrac{AF}{AC}=\dfrac{1}{4}\cdot\dfrac{3}{4}=\dfrac{3}{16},
\]
and in the same way $[BED]=\dfrac{3}{4}\cdot\dfrac{1}{4}[ABC]$ and $[CFE]=\dfrac{3}{4}\cdot\dfrac{1}{4}[ABC]$. Removing these three corner triangles,
\[
[DEF]=\left(1-3\cdot\dfrac{3}{16}\right)\cdot48=\dfrac{7}{16}\cdot48=\boxed{21} .
\]
